\documentclass{article}

\begin{document}
    \begin{center}
        \Large \textbf{Challenge 3: DNS Records und Firewall} \\ [6pt]
        \normalsize
        Autor: Maximilian Jakob Maag B.Sc. \\
        Version: 1.0
    \end{center}
    \section{Intro}
        Ein VPS allein genügt oft nicht. Es müssen viele Einstellungen getroffen werden, die nicht auf dem VPS selbst festgelegt werden können.
        Für einen Mailserver müssen z.B. A-Records, MX-Records und diverse TXT Records angelegt werden. \\ \\
        Diese Übung setzt voraus, dass Sie einen Mailserver deployen können. \\ \\
        U.a. der Provider Linode lässt das setzen von DNS Records und andere Einstellungen einer Domain per Terraform zu. \\ \\
        Neben DNS Records spielen Firewalls eine wichtige Rolle für die IT Security. Das Erstellen einer Firewall, das Zuweisen von VPSs und das Erstellen der
        entsprechenden Firewallregeln kann ebenfalls mit Terraform vorgenommen werden. Je nach Provider unterscheiden sich die unterstützten Features.
    \section{Aufgaben}
        In dieser Challenge ist es Ihr Ziel, passend zu Ihrem Mailserver dessen DNS-Records mittels Terraform zu setzen und
        eine ausreichend schützende Firewall zu konfigurieren.
    \subsection{Aufgabe 1}
        Testen Sie Ihre Konfiguration aus den letzten Challenges und deployen Sie einen Mailserver, z.B. Mailcow. \\
        Anschließend zerstören Sie diesen wieder, damit er nicht ungeschützt im Internet erreichbar ist.
    
    \subsection{Aufgabe 2}
        Erstellen Sie die Konfiguration für eine Firewall. Blockieren Sie zunächst sämtlichen Inbound- und Outbound-Traffic.
        Fügen Sie Ihren Mailserver und mindestens drei andere VPS der Firewall hinzu. Pingen Sie die Instanzen an. Was beobachten Sie?

    \subsection{Aufgabe 3}
        Ändern Sie die Konfiguration Ihrer Firewall ab. Lassen Sie ping und ssh zu.

    \subsection{Aufgabe 4}
        Ändern Sie die Konfiguration Ihrer Firewall erneut ab. Erarbeiten Sie sinnvolle Einstellungen für den Betrieb ihres Mailservers.
        Testen Sie ob alles funktioniert, indem Sie auf Ihrem deployten Mailserver ein E-Mailkonto einrichten und eine E-Mail an sich selbst versenden.

    \subsection{Aufgabe 5}
        Verbinden Sie sich mit einem der anderen Beispiel VPS aus Aufgabe 1. Installieren Sie dort einen Webserver, der eine simple statische Webseite bereitstellt.
        Ändern Sie die Regeln Ihrer Firewall. An diesem Server darf kein Mailport offen sein. Nur die Ports 443 und 80 für den Webserver, sowie 22 für SSH dürfen offen sein.
        Der Mailserver muss von dieser Änderung unberührt bleiben und weiterhin E-Mails versenden können.
    
    \subsection{Aufgabe 6}
        Legen Sie einen A‑ bzw. einen AAAA‑Record fest, der auf Ihren Webserver zeigt. Rufen Sie Ihre Webseite auf. Der Browser muss Ihnen eine statische Webseite ausliefern mittels einer gültigen HTTPS‑Verbindung.
        Führen Sie am Ende der Übung terraform destroy aus.

\end{document}